\chapter{Traitement Multilingue: Pipeline Retenu et Variante}

\section{Pipeline retenu pour l'evaluation finale}
Le pipeline principal ne fusionne pas explicitement les faits EN et FR dans une
seule timeline d'entrainement. Il procede ainsi:
\begin{itemize}[leftmargin=2em]
  \item extraction d'une timeline \texttt{en},
  \item extraction d'une timeline \texttt{fr},
  \item traduction de \texttt{fr} vers \texttt{fr\_en} via LLM,
  \item apprentissage de similarite sur des paires \texttt{en <-> fr\_en}.
\end{itemize}

Ce choix donne un signal supervise cross-lingue direct pour la tache de
similarite narrative, tout en evitant une supervision redondante liee a une
fusion trop proche du texte EN source.

\section{Role du LLM dans cette etape}
Dans le pipeline retenu, le LLM est utilise pour la traduction FR->EN
(\texttt{scripts/translate\_fr\_timelines\_to\_en.py}).
La similarite est ensuite apprise par sentence-encoder fine-tune
(allMiniLM, MPNet).

\section{Variante exploree: fusion semantique + final\_en}
Une variante a ete implementee:
\begin{itemize}[leftmargin=2em]
  \item fusion semantique FR/EN (\texttt{scripts/merge\_multilingual\_timelines.py}),
  \item finalisation monolingue EN unique (\texttt{scripts/finalize\_monolingual\_en\_timelines.py}).
\end{itemize}

Cette variante produit des fichiers \texttt{*.final\_en.timeline.json}
(249 entites), utiles pour l'analyse, mais les scores finaux presentes dans ce
rapport sont calcules sur la branche \texttt{en/fr\_en}.
